For $x>a>0$, every monomial in $q(x)$ is nonnegative, so $(x-a)q(x)\ge0$ and $f(x)\ge r>0$. Thus no such $x$ is a root.
Checking only $f(a)>0$ fails: $f(x)=(x-11)(x-12)$ has $f(10)=2>0$ but has two roots greater than $10$.

For the given cubic, synthetic division at $a=4$ gives
\[
x^3-2x^2-3x-4=(x-4)(x^2+2x+5)+16.
\]
The criterion therefore proves that every real root is at most $4$. In fact $f(4)=16$, so none equals $4$. Since $f(3)=-4$, continuity also gives a root in $(3,4)$. The quotient coefficients and the remainder, not just the single value at the test point, are what justify the upper bound.
